\section{Methods}

\subsection{System Preparation}

Initial coordinates were taken from the mouse Cys-GPX6 crystal structure (PDB~7FC2), and an AlphaFold~2-predicted model was built for human Sec-GPX6 (UniProtKB P59796).\cite{PDB7FC2,jumper2021alphafold} 
Crystallographic solvent and nonessential hydrogen atoms were removed before setup. 

Residue labels throughout the study use the renumbered coordinate model prepared from 7FC2. The PDB entry retains UniProt numbering but reports residues 1--24 as unresolved; these residues were omitted and the retained coordinates, beginning at PDB residue~25, were renumbered from~1 for simulation setup. Study residues~49, 83, and 157 therefore correspond to PDB/UniProt residues~73, 107, and 181, respectively. This study numbering is used consistently in the manuscript, figures, Supporting Information, and input files.

Protein sequences for human GPX6 (UniProtKB P59796) and mouse GPX6 (UniProtKB Q91WR8) were retrieved from UniProt and aligned with Clustal Omega 1.2.4\cite{Sievers2011ClustalOmega} to identify substitutions between the two proteins (Figure~\ref{SI-fig:gpx6_sequence_alignment}). We selected the 20 positions within 20~\AA{} of the catalytic chalcogen; the substitutions and distances in both directions are reported in Tables~\ref{SI-tab:mouse_mutation_distance} and \ref{SI-tab:human_mutation_distance}. This distance includes residues in the first coordination shell as well as more distant positions that may affect the structure or electrostatics of the active site, while keeping the iterative EVB calculations tractable. The selected region lies within the 25~\AA{} SCAAS solvent sphere, but its boundary was chosen independently.

\subsection{EVB Model and Parameterization}

We used the standard two-state free-energy perturbation/empirical valence bond (FEP/EVB) formulation, in which reactant and product diabatic potentials are sampled through intermediate mapping windows and subsequently mixed to reconstruct the EVB ground-state free-energy profile. The Hamiltonian, mapping potentials, energy-gap coordinate, FEP accumulation, and reweighting equations are provided in the Supporting Information, together with the complete workflow (Figure~\ref{SI-fig:workflow}).\cite{warshel_empirical_1980,marelius_q_1999,oanca_efficient_2023}

The reactive region comprised the catalytic Cys/Sec residue at position~49, hydrogen peroxide, and Gln83, which together define the proton-transfer step that generates the catalytic thiolate or selenolate. Partial charges for the reactive atoms in the different diabatic states were derived using Schrödinger's \texttt{ffld\_server} tool in Maestro 14.5 with the optimized potentials for liquid simulations--all atom (OPLS-AA) force field.\cite{Schrodinger2024Maestro} State-specific parameters for the EVB reactive region, including Morse terms and reactant/product charges for the sulfur-, selenium-, peroxide-, and glutamine-containing states, are summarized in Tables~\ref{SI-tab:morse} and \ref{SI-tab:partial_charges}; the reactive-atom naming is provided in the Supporting Information (Figure~\ref{SI-fig:atoms}). In the Q FEP files, Morse bond terms were used only for bonds whose connectivity changes between the two diabatic states: the peroxide O1--H1 and chalcogen--HG1 bonds in the reactant state, and the O1--HG1 and Gln83 OE1--H1 bonds in the product state. These Morse terms are part of the EVB state potentials sampled by the mapping windows. Other local bond adjustments, such as changes within Gln83 and the Sec C$\beta$--Se bond, were described with harmonic bond terms in the FEP topology and should not be confused with the external flat-bottom distance restraints described below.

The EVB state-energy shift $\alpha$ and coupling $H_{12}$ were calibrated against the quantum-mechanical free-energy profile reported for the human Sec-containing glutathione peroxidase mechanism,\cite{prabhakar_is_2006} using $\Delta G^{\ddagger}=16.40$~kcal/mol and $\Delta G^{\circ}=-60.52$~kcal/mol as reference targets. $H_{ij}$ can be expressed as a distance-dependent off-diagonal function or as a function of other variables, such as the energy gap between the two EVB states, but in this work we used a constant coupling. The resulting calibrated parameter set (Table~\ref{SI-tab:evb_parameters}) was then kept fixed across all mutant and homolog calculations so that energetic differences reflect changes in the protein scaffold rather than re-fitting of the EVB surface.

All calculations were done using Q6.\cite{marelius_q_1999} Input generation and EVB analysis used qtools, Qprep, Qdyn, Q free-energy perturbation (QFEP), and QMapper; the nonreactive region used OPLS-AA.\cite{PurgBauer2017Qtools,jorgensen1996opls_aa} In the Qdyn FEP simulations the predefined off-diagonal term was set to zero, so trajectories sampled the mapping potentials. QMapper then applied EVB mixing and the calibrated state-energy shift during analysis. These implementation choices, which are specific to the present calculations, are detailed in the Supporting Information.

\subsection{Geometric Restraints, Solvation, and Relaxation}

The reactive region was placed in a 25~\AA{} sphere of water centered on the catalytic sulfur or selenium atom of residue~49. Spherical boundary conditions were treated with the surface constrained all-atom solvent (SCAAS) model, which constrains the finite solvent boundary so that it approximates a bulk environment.\cite{king_surface_1989} All simulations used 10~\AA{} cutoffs for solute--solute, solute--solvent, and solvent--solvent interactions, long-range electrostatic corrections, and the SHAKE bond-constraint algorithm for solvent bonds.

The initial structure-relaxation inputs were generated with the qtools script \texttt{q\_genrelax.py}\cite{PurgBauer2017Qtools} using the EVB FEP file, but with the Q lambda vector fixed at \texttt{1.00 0.00}, corresponding to the reactant diabatic state. This initial relaxation did not include the active-site flat-bottom distance restraints used for FEP sampling. The template comprised two 25.1~ps stages at low temperature: solvent relaxation with sequence restraints on the solute, followed by relaxation without sequence restraints, yielding 50.2~ps of preproduction relaxation. The minimized restart from this step was used to write the structure for subsequent EVB sampling.

Flat-bottom harmonic distance restraints were introduced during the pre-FEP equilibration and retained during the production FEP windows to preserve a chemically productive reactive geometry while still allowing local motion along the proton-transfer coordinate.  The restrained coordinates were the chalcogen--peroxide O1 and Gln83 OE1--peroxide O1 donor--acceptor distances, the reactant chalcogen--HG1 and peroxide O1--H1 bonds, and the two product-side proton-transfer distances Gln83 OE1--H1 and peroxide O1--HG1 (Table~\ref{SI-tab:flatbottom}); the corresponding potentials are provided in the Supporting Information (Figure~\ref{SI-fig:flatbottom}). Before each EVB production profile, systems were heated for 2~ps at 1~K, 3~ps at 100~K, and 5~ps at 300~K, followed by six 10~ps restraint-release stages, giving 70~ps of pre-FEP equilibration. During this pre-FEP stage, sequence restraints were progressively reduced and then removed, while the active-site flat-bottom restraints were reduced from their initial heat-up values to the final production value of 2~kcal~mol$^{-1}$~\AA$^{-2}$, which was retained in all production FEP windows. The stepwise protocol is summarized in Table~\ref{SI-tab:combined_protocol}.

\subsection{EVB Free-Energy Calculations}

Production EVB profiles were obtained from 51 FEP windows generated from $\boldsymbol{\lambda}=(1.0,0.0)$ to $\boldsymbol{\lambda}=(0.0,1.0)$, i.e., from the reactant mapping potential to the product mapping potential. The production template used a 1.0~fs timestep at 300~K. Each FEP window was sampled for 20\,000 steps (20~ps), corresponding to 1.02~ns per complete EVB profile. Energy files were written every 5 molecular dynamics (MD) steps, giving 4000 saved energy records per window and 204\,000 records per complete profile per replica before QFEP point skipping. Free-energy surfaces were reconstructed with 50 energy-gap bins, requiring at least 10 points per bin and excluding the first 100 sampled points from the final analysis. Activation free energies ($\Delta G^{\ddagger}$) and reaction free energies ($\Delta G^{\circ}$) were obtained from the bin-averaged, reweighted EVB profiles, and uncertainties are reported as the standard error across replicas. QFEP could still return the end-to-end FEP mapping value $\Delta G_{\lambda}$ when the adjacent-window integration was complete but the energy-gap data did not satisfy the bin-occupancy requirements for reconstructing a complete EVB profile. Consequently, Table~\ref{SI-tab:evb_stats} reports a separate $N$ for $\Delta G_{\lambda}$; these additional mapping values were not used to calculate $\Delta G^{\ddagger}$ or $\Delta G^{\circ}$.

Each complete 51-window EVB profile was treated as one statistical observation. Individual MD records and energy-gap bins were not counted as independent observations. $N$ is therefore the number of independently initiated profiles that met the same criteria for successful profile reconstruction. We calculated the standard error of the mean (SEM) from the variation among profiles as $\mathrm{SEM}=\mathrm{SD}/\sqrt{N}$, where SD is the sample standard deviation. We did not use autocorrelation estimates from individual windows to rescale these uncertainties. Each biased window is only 20~ps long, and its autocorrelation describes local fluctuations in the energy or energy gap, not the uncertainty of the barrier reconstructed from all 51 windows. Moreover, there is no instantaneous barrier value for each saved frame. The independent 1.02~ns profiles are therefore the relevant repeats for estimating uncertainty in the free energies. Oanca and \AA{}qvist analyzed the robustness of EVB-derived activation parameters in computational Arrhenius plots; the present calculations were performed at a single temperature and do not constitute an Arrhenius analysis.\cite{oanca_arrhenius_2024}
Representative active-site configurations from the reactant, transition-state-region, and product mapping windows are shown in Figure~\ref{SI-fig:evb_active_site_configurations}.

\subsection{Mutational Pathway Analysis}

Human and mouse GPX6 differ at 47 positions. We focused on the 20 residues within 20~\AA{} of the catalytic chalcogen, using the rationale described above. Each variant was prepared with the same structural and EVB protocol used for the residue-49 exchange. We first replaced Sec49 with Cys in human GPX6 and Cys49 with Sec in mouse GPX6. These two variants define level~0 of the human-to-mouse and mouse-to-human paths, respectively. At each subsequent level, we calculated all remaining single substitutions and retained the one with the lowest activation barrier before proceeding to the next level. This procedure yields one low-barrier path in each direction; it does not imply that either path is unique.

\subsection{Phylogenetic and Sequence-Design Analyses}

To place the EVB-derived trajectories in phylogenetic context, the GPX6 maximum-likelihood framework and reconstructed ancestral sequence were taken from Rees et al.\cite{rees_ancient_2024} The ancestral reference used here is the reconstructed Sec-containing ancestor subtending the primate and rodent/lagomorph branches, referred to as node~25 in that study. Pairwise sequence distances among the EVB-generated human and mouse GPX6 variants and this ancestral sequence were calculated over the focused 20-residue subset used in the pathway analysis as normalized sequence distances, $d_{ij}=1-I_{ij}$, where $I_{ij}$ is the fraction of identical residues in the compared subset. These distances were visualized as a pairwise distance heatmap with Ward hierarchical clustering and projected into a triangular geometric representation of sequence space.

We used the Protein message-passing neural network (ProteinMPNN) to test whether the structures sampled along the EVB paths favor sequences similar to the reconstructed ancestor.\cite{Dauparas2022ProteinMPNN} We generated sequence ensembles for all available endpoint and intermediate structures from both paths. For each level, the input was the structure sampled in the $\lambda=0.5$ EVB window, which represents the transition-state region of the profile. ProteinMPNN was run with the soluble-protein model and a fixed backbone, generating 100 sequences for each structure.

To visualize the positioning of sequences relative to the modern endpoints and the ancestor, we projected all sequences into a two-dimensional triangular coordinate system. The three vertices of the equilateral triangle were defined by the mouse GPX6 wild-type sequence at the bottom-left, the human GPX6 wild-type sequence at the bottom-right, and the node~25 ancestral sequence at the apex. For every query sequence, we computed three pairwise distances to the reference sequences using the normalized Levenshtein metric,\cite{LiLiu2007NormalizedLevenshtein} defined here as the edit distance divided by the longer of the two compared sequence lengths. Distances to the three vertices were converted into barycentric weights using $w_i=\exp(-d_i/\tau)/\sum_j\exp(-d_j/\tau)$, with a dimensionless softmax temperature $\tau=1.0$ for both the EVB-pathway and ProteinMPNN projections. The final Cartesian coordinate was the weighted sum of the three triangle vertices. This triangular representation enabled direct comparison of how the greedy EVB trajectories and the ProteinMPNN designs explore sequence space and approach the reconstructed ancestor.
